% !TEX root = ../main.tex
% ============================================================================
% Owner: Arteaga
%
% Target length: at most 1.5 pages. Each paper gets what it does and the one
% fact we use; the contrast with our work is stated once, in the closer.
%
% Papers in Docs/Papers/ are named by arXiv id, one PDF per bib entry.
% Every bib entry was checked against arXiv on 2026-11-06.
% ============================================================================
\section{Literature Review}
\label{sec:lit}

\paragraph{Monocular depth foundation models.}
Our backbone belongs to a family of models trained to predict relative depth from a single image across many datasets. MiDaS \citep{ranftl2020midas} trains on datasets with incompatible annotations by using a loss that aligns prediction and ground truth in scale and shift before comparing them, so the network predicts inverse depth only up to those two unknowns, and they are fitted per image by least squares before any error is measured. Our per-image fits against lidar and radar are the same alignment, done with sparse sensor returns. DPT \citep{ranftl2021dpt} replaces the convolutional encoder with a vision transformer whose tokens are reassembled into feature maps at four resolutions and merged by a convolutional decoder, followed by a short convolutional head. Depth Anything \citep{yang2024depthanything} and Depth Anything V2 \citep{yang2024depthanythingv2} pair a DINOv2 encoder with this decoder and scale the training data by pseudo-labeling about 62M unlabeled images with a teacher model. The first convolution of the head, named \texttt{output\_conv1} in the V2 code, gives the 64-channel map for ViT-B that our head reads. The released models output affine-invariant inverse depth and no per-pixel confidence; we use the relative ViT-B checkpoint frozen (deviation D1 in Table~\ref{tab:deviations}).

\paragraph{Uncertainty and confidence for monocular depth.}
Most work on uncertainty for depth changes the depth model. \citet{kendall2017uncertainties} separate aleatoric uncertainty, the noise inherent in the observations, from epistemic uncertainty, which lies in the model, and estimate the first with a variance output learned through the loss and the second with Monte Carlo dropout \citep{gal2016dropout}, which runs the network several times with resampled dropout masks. Deep ensembles \citep{lakshminarayanan2017ensembles} train several networks from different initializations and combine their predictions. All three are trained with the depth network or need several copies of it, which a single frozen checkpoint rules out. Our proposal listed MC dropout as a baseline; we drop it because every dropout rate in the Depth Anything V2 encoder is zero and the decoder has none, so repeated passes are identical (Section~\ref{sec:progress}).

\citet{poggi2020uncertainty} compare eleven ways of estimating per-pixel uncertainty for a self-supervised depth network (Monodepth2 on KITTI) and supply our evaluation protocol. Pixels are removed in order of decreasing uncertainty and the error of the remainder is plotted. AUSE, introduced for optical flow by \citet{ilg2018uncertainty}, is the area between that sparsification curve and the oracle curve that removes pixels by true error (lower is better); AURG is the area between the flat curve of a random ordering and the model's curve (higher is better). The paper gives both in words, so how the areas are computed over our labeled grid cells is our own choice. Its ``Post'' method, the difference between the depth of an image and the flipped-back depth of its mirror image, needs no retraining, and our flip-disagreement baseline is adapted from it. The paper also notes that a classifier trained on the prediction error, which is what our head is, is common in stereo and optical flow but needs ground-truth depth; our lidar supplies it.

\citet{hornauer2022gradient} estimate uncertainty for a trained depth network without changing its weights and without any training: the difference between the depth maps of an image and its horizontal flip is used as a loss, and the gradient of that loss at one decoder layer, pooled over channels, is the uncertainty map. On Monodepth2 it ranks errors better than the plain flip difference, at the cost of a backward pass per image. It is among the nearest prior work to ours in leaving the depth model untouched and reading its decoder, but it uses no ground truth.

\paragraph{Radar-camera fusion.}
Recent work fuses radar with images to improve depth or 3D perception; we instead give radar to the head as evidence of where the frozen depth disagrees with an independent range measurement. \citet{lin2020radardepth} predict dense depth on nuScenes from an image and sparse projected radar, and list what makes radar harder to use than lidar: far fewer points, a limited vertical field of view, noise from multipath reflections, and disagreement with lidar even for clean returns. Their network discards radar points that disagree with a first depth estimate. Our radar residual is the same comparison, passed to the head as an input instead of being thresholded, and its noise in our data (Section~\ref{sec:results}) agrees with their account. HyDRa \citep{wolters2025hydra} and LiCROcc \citep{ma2024licrocc} bring radar into networks for 3D detection and occupancy, the first by fusing radar features before the depth network and the second by distilling a lidar-camera teacher into a radar-camera student; Occ3D \citep{tian2023occ3d} is the occupancy benchmark for our stretch goal (Section~\ref{sec:plan}). BEVCar \citep{schramm2024bevcar} fuses cameras and radar for bird's-eye-view segmentation and releases a division of the nuScenes validation set into day, rain, and night groups of 4,449, 968, and 602 samples, which we use for our per-condition results.

\paragraph{Adverse-condition driving datasets.}
nuScenes \citep{caesar2020nuscenes}, the dataset we train and evaluate on, provides 1,000 driving scenes of 20 seconds with six cameras, a 32-beam lidar, and five radars, with rain and night at 19.4\% and 11.6\% of keyframes. Camera exposure is triggered as the lidar sweeps across each camera's field of view, which lets us project lidar into the image as labels, and the radar sweeps between keyframes are released, which our radar residual accumulates. md4all \citep{gasperini2023md4all} makes a depth network robust at night and in rain by training on day images translated to those conditions; evaluated per condition against lidar on nuScenes, its self-supervised baselines show depth error roughly doubling from day to night, the direction we see in our own cache (Section~\ref{sec:results}).
% Access terms ("by request") are not in the paper; cite the project page before stating them here.
Seeing Through Fog \citep{bijelic2020seeing}, STF below, has the fog and snow that nuScenes lacks, with lidar and radar, and is our candidate for a second evaluation, though adverse frames are a small part of its labeled set and lidar's perceivable distance falls to about 25\,m in fog.
% The radar added in the 2025 ZOD release is not in this paper; cite the release before mentioning it.
The Zenseact Open Dataset \citep{alibeigi2023zod}, ZOD below, offers 20-second camera and lidar sequences but, as described in the paper, no radar, so it is a fallback (Section~\ref{sec:plan}).

\paragraph{Where this work sits.}
The depth foundation models return relative depth with no indication of reliability; we take one as given and add that indication. The uncertainty literature mostly builds the estimate into the depth network or samples it, which a frozen public checkpoint does not allow, and the post-hoc methods that remain use no ground truth; we train a small head on lidar-derived error labels and evaluate it with that literature's sparsification measures. Radar-camera work uses radar to improve the prediction itself; we use it to predict where the depth is wrong. Whether the frozen decoder features carry enough information to predict the model's own error, and how much radar and temporal consistency add, are the questions the experiments answer.
